% =============================================================================
%  SPECTRA — Paper Section Draft
%  Label Provenance & Evaluation Methodology
%  Target: IEEE Transactions on Information Forensics & Security
%  Author: Sagar Korde
% =============================================================================
%
%  INSERT this block inside Section V (Experimental Evaluation),
%  as a subsection before Table I (Baseline Comparison).
% =============================================================================

% ─── V-B  Label Provenance and Evaluation Methodology ────────────────────────

\subsection{Label Provenance and Evaluation Methodology}
\label{subsec:label_provenance}

\subsubsection{Dataset and Transaction Type Labels}
We evaluate on the CryptoGraphNet dataset~\cite{cryptographnet2024}, which contains
\num{300000} stratified Bitcoin transactions sampled across six mutually exclusive
transaction types: Standard, Peer-to-Peer (P2P), Consolidation, Distribution,
Batch-Payment, and CoinJoin-like.
The six class labels are assigned via deterministic heuristics applied to
structural transaction attributes---specifically, the number of distinct input
and output addresses, value concentration ratios, and fee-rate patterns.
For example, a \emph{Consolidation} transaction is defined as one with
$|\mathit{inputs}| \geq 3$ and $|\mathit{outputs}| = 1$, while a
\emph{Batch-Payment} requires $|\mathit{outputs}| \geq 5$~\cite{cryptographnet2024}.

\subsubsection{Label--Feature Entanglement}
This construction introduces a structural entanglement between labels and tabular
features: because the class membership rules are algebraic functions of
$|\mathit{inputs}|$, $|\mathit{outputs}|$, and value statistics, any model
that observes these raw tabular features can reconstruct the labels
with near-perfect accuracy without learning any generalizable behavioral
pattern.
Table~\ref{tab:fair_baselines} confirms this: XGBoost and
BERT4ETH-BTC both achieve $\text{F}_1 = 1.000$ on the original feature set.
Critically, this degeneracy \emph{persists even after temporal identifiers
are removed} (Table~\ref{tab:fair_baselines}, ``fair'' columns), because
the definitive features are transaction-structural, not temporal.

We emphasize that this is a property of the benchmark's label-generation
procedure---a known limitation of rule-based annotation common to
blockchain datasets~\cite{Jourdan2019,Akcora2022}---and not an indication
of model overfitting.
We report tabular-classifier results for completeness per IEEE TIFS
reproducibility requirements, and clearly distinguish them from the
more informative \emph{graph-topology} track discussed below.

\subsubsection{Two-Track Evaluation Protocol}
To provide a rigorous and interpretable comparison, we partition all methods
into two evaluation tracks:

\begin{enumerate}[leftmargin=*, label=\textbf{T\arabic*}]

  \item \textbf{Tabular-feature classifiers} (XGBoost, BERT4ETH-BTC).
    These methods receive the full per-transaction feature vector
    $\mathbf{x} \in \mathbb{R}^{20}$ selected by Mutual Information ranking.
    Because $\mathbf{x}$ includes $|\mathit{inputs}|$, $|\mathit{outputs}|$,
    and value statistics from which the labels are derived, their accuracy
    constitutes an \emph{upper-bound oracle} rather than a practically
    meaningful measure of fraud-detection capability.

  \item \textbf{Graph-topology methods} (SPECTRA, EvolveGCN, AA-GNN,
    clustering baselines).
    These methods operate exclusively on the address-level transaction graph
    $\mathcal{G} = (\mathcal{V}, \mathcal{E}, \mathbf{W})$ and/or the
    spectral node embedding $\boldsymbol{\Phi} \in \mathbb{R}^{|\mathcal{V}| \times k}$
    derived from the graph Laplacian.
    Raw structural features ($|\mathit{inputs}|$, $|\mathit{outputs}|$) are
    \emph{not} directly available to these methods; the graph topology
    encodes them only implicitly through node degree distributions.
    This track constitutes the \emph{meaningful} comparison and is the
    focus of SPECTRA's contribution claims.

\end{enumerate}

\subsubsection{Temporally-Fair Feature Set}
Independent of the two-track split, we additionally ablate temporal
identifier features.
In Bitcoin's ledger, \texttt{block\_height} and \texttt{block\_time}
are monotonically increasing proxies for calendar time.
Because different transaction types were prevalent at different epochs
of Bitcoin's adoption curve (e.g., CoinJoin-like transactions became
common only after 2017, and batch payments correlate with exchange activity
peaks in 2020--2021), a model with access to these fields can exploit
distributional shift rather than structural features.
We therefore define a \emph{temporally-fair} feature set
$\mathbf{x}^{\dagger}$ by removing \texttt{block\_height},
\texttt{block\_time}, and \texttt{year}, then re-selecting the top-20
features by Mutual Information on the remaining columns.
All results are reported for both the original set $\mathbf{x}$ and
the fair set $\mathbf{x}^{\dagger}$ in Table~\ref{tab:combined_comparison}.
Cyclical time features (\texttt{hour}, \texttt{day\_of\_week},
\texttt{week\_of\_year}) are retained, as they reflect genuine
user-activity rhythms~\cite{Ron2013} rather than temporal identity.

\subsubsection{Why SPECTRA Is Inherently Temporally Fair}
SPECTRA's GNN component (Module~R) receives node features
$\mathbf{h}_v \in \mathbb{R}^{d}$ constructed from:
(i) spectral graph embeddings $\boldsymbol{\phi}_v$ (Laplacian eigenvectors),
(ii) address-level behavioral statistics (fan-in ratio, transaction velocity,
average received/sent value, inter-transaction gap),
(iii) VAE-derived latent coordinates $\mathbf{z}_v$,
(iv) Bayesian trust score $\tau_v$, and
(v) Markov chain anomaly score $\gamma_v$.
None of these quantities encodes absolute blockchain time or block height.
SPECTRA's classification results therefore belong unambiguously to Track~T2
and are directly comparable with EvolveGCN and AA-GNN without any
temporal-fairness correction.

% ─── References used above (add to bibliography) ─────────────────────────────
%
%  \bibitem{cryptographnet2024}
%    Author(s), ``CryptoGraphNet: A Labeled Bitcoin Transaction Graph Dataset,''
%    [Online]. Available: [URL]. 2024.
%
%  \bibitem{Jourdan2019}
%    M.~Jourdan, S.~Blandin, L.~Wynter, and P.~Deshpande,
%    ``Characterizing entities in the Bitcoin blockchain,''
%    in \textit{Proc. IEEE Int. Conf. Data Mining Workshops (ICDMW)}, 2019.
%
%  \bibitem{Akcora2022}
%    C.~G.~Akcora, Y.~Li, Y.~R.~Gel, and M.~Kantarcioglu,
%    ``BitcoinHeist: Topological Data Analysis for Ransomware Prediction
%    on the Bitcoin Blockchain,''
%    in \textit{Proc. IJCAI}, 2022, pp.~4439--4445.
%
%  \bibitem{Ron2013}
%    D.~Ron and A.~Shamir,
%    ``Quantitative Analysis of the Full Bitcoin Transaction Graph,''
%    in \textit{Proc. Financial Cryptography and Data Security (FC)}, 2013,
%    pp.~6--24.
